% GENERATED FILE. Edit canonical lessons, then run npm run generate:books.
% canonical-source-hash: fnv1a64:f7dbe760a7e0bfb7
% canonical-lessons: PA-C190-karna-return, PA-C190-kita, PA-C190-kita-review, PA-C190-kamm-kita, PA-C190-frame-review, PA-C190-infinitive-perfective-contrast, PA-C190-payoff

\chapter{I Did the Work}
\label{ch:pa-a2-kita-perfective}
\hlchaptermodality{190}

\begin{chapteropening}
\textbf{By the end of this chapter:} \emph{I can distinguish the familiar infinitive from one bounded masculine-singular completed-action frame.}

Retrieve the infinitive and one masculine-singular completed-action frame through listening, speaking, reading, and writing.
\end{chapteropening}

\section[\texorpdfstring{\pa{ਕਰਨਾ} · \pa{ਮੈਂ} \pa{ਅਮਨ} \pa{ਨੂੰ} … (karnaa; main …)}{ਕਰਨਾ · ਮੈਂ ਅਮਨ ਨੂੰ … (karnaa; main …)}]{Return to do and a completed action}

\label{lesson:PA-C190-karna-return}



\begin{quote}
\emph{Hear:} \emph{karnaa}; YOU SAY: \textbf{\pa{ਕਰਨਾ}}
\end{quote}

\begin{grammarlens}[title={Grammar Lens}]
Retrieve the familiar completed-action frame before changing the verb: \textbf{\pa{ਮੈਂ} \pa{ਅਮਨ} \pa{ਨੂੰ} \pa{ਮਿਲਿਆ}}.
\end{grammarlens}

\subsection*{Your turn}
\begin{itemize}
\raggedright
  \item \emph{Hear:} \emph{karnaa}; YOU SAY: \textbf{\pa{ਕਰਨਾ}}
  \item \emph{Hear:} \emph{main Aman nuun miliaa}; YOU SAY: \textbf{\pa{ਮੈਂ} \pa{ਅਮਨ} \pa{ਨੂੰ} \pa{ਮਿਲਿਆ}}
\end{itemize}

\begin{tcolorbox}[breakable,colback=teal!4,colframe=teal!35!black,title={Before you move on}]
Say the infinitive, then the familiar completed-action frame. This opener owns no new form.
\end{tcolorbox}
\section[\texorpdfstring{\pa{ਕੀਤਾ} (kiitaa)}{ਕੀਤਾ (kiitaa)}]{The completed form of \pa{ਕਰਨਾ}}

\label{lesson:PA-C190-kita}



\begin{quote}
Retrieve \textbf{\pa{ਕਰਨਾ}}, then the final word in \textbf{\pa{ਮੈਂ} \pa{ਅਮਨ} \pa{ਨੂੰ} \pa{ਮਿਲਿਆ}}.
\end{quote}

\begin{grammarlens}[title={Grammar Lens: direct owner}]
\begin{quote}
\textbf{\pa{ਕੀਤਾ}} — \emph{kiitaa} — did or made, masculine singular completed form
\end{quote}

Treat the whole written form as today's measured target. The next lesson retrieves it before a complete sentence frame is built; the later frame will make its agreement controller explicit.
\end{grammarlens}

\subsection*{Your turn}
\begin{itemize}
\raggedright
  \item \emph{Hear:} \emph{kiitaa}; YOU SAY: \textbf{\pa{ਕੀਤਾ}}
  \item \emph{Read it:} \textbf{\pa{ਕੀਤਾ}}
  \item \emph{Write it:} copy \textbf{\pa{ਕੀਤਾ}} once
\end{itemize}

\begin{tcolorbox}[breakable,colback=teal!4,colframe=teal!35!black,title={Before you move on}]
Produce the completed form without looking.
\end{tcolorbox}
\section[\texorpdfstring{\pa{ਕੀਤਾ} (kiitaa)}{ਕੀਤਾ (kiitaa)}]{Retrieve the completed form}

\label{lesson:PA-C190-kita-review}



\begin{quote}
Hear “did or made” and prepare the single completed form.
\end{quote}

\begin{grammarlens}[title={Grammar Lens}]
Keep \textbf{\pa{ਕੀਤਾ}} distinct from the infinitive \textbf{\pa{ਕਰਨਾ}}. Familiar \textbf{\pa{ਮਿਲਿਆ}} is a completed-action comparison, not an agreement template. In the coming transitive frame, \textbf{\pa{ਕੀਤਾ}} will agree with the unmarked masculine-singular object \textbf{\pa{ਕੰਮ}}.
\end{grammarlens}

\subsection*{Your turn}
\begin{itemize}
\raggedright
  \item \emph{Hear:} \emph{kiitaa}; YOU SAY: \textbf{\pa{ਕੀਤਾ}}
  \item \emph{Read it:} \textbf{\pa{ਕੀਤਾ}}
  \item \emph{Write it:} cover the model and write \textbf{\pa{ਕੀਤਾ}}
\end{itemize}

\begin{tcolorbox}[breakable,colback=teal!4,colframe=teal!35!black,title={Before you move on}]
Read the form once more and state its bounded job.
\end{tcolorbox}
\section[\texorpdfstring{\pa{ਮੈਂ} \pa{ਕੰਮ} \pa{ਕੀਤਾ} (main kamm kiitaa)}{ਮੈਂ ਕੰਮ ਕੀਤਾ (main kamm kiitaa)}]{Build one completed-action frame}

\label{lesson:PA-C190-kamm-kita}



\begin{quote}
Retrieve \textbf{\pa{ਮੈਂ}}, \textbf{\pa{ਕੰਮ}}, and \textbf{\pa{ਕੀਤਾ}}.
\end{quote}

\begin{grammarlens}[title={Grammar Lens: complete statement}]
\begin{quote}
\textbf{\pa{ਮੈਂ} \pa{ਕੰਮ} \pa{ਕੀਤਾ}} — \emph{main kamm kiitaa} — I did the work
\end{quote}

Read the familiar speaker, familiar work noun, and completed form as one bounded frame. In this transitive perfective frame, \textbf{\pa{ਕੀਤਾ}} agrees with the unmarked masculine-singular object \textbf{\pa{ਕੰਮ}}, not the speaker.
\end{grammarlens}

\subsection*{Your turn}
\begin{itemize}
\raggedright
  \item \emph{Hear:} \emph{main kamm kiitaa}; YOU SAY: \textbf{\pa{ਮੈਂ} \pa{ਕੰਮ} \pa{ਕੀਤਾ}}
  \item \emph{Read it:} \textbf{\pa{ਮੈਂ} \pa{ਕੰਮ} \pa{ਕੀਤਾ}}
  \item \emph{Write it:} copy \textbf{\pa{ਕੰਮ} \pa{ਕੀਤਾ}} once
\end{itemize}

\begin{tcolorbox}[breakable,colback=teal!4,colframe=teal!35!black,title={Before you move on}]
Produce the complete “I did the work” frame.
\end{tcolorbox}
\section[\texorpdfstring{\pa{ਮੈਂ} \pa{ਕੰਮ} \pa{ਕੀਤਾ} (main kamm kiitaa)}{ਮੈਂ ਕੰਮ ਕੀਤਾ (main kamm kiitaa)}]{Retrieve the completed work frame}

\label{lesson:PA-C190-frame-review}



\begin{quote}
Hear “I did the work” and prepare the whole statement.
\end{quote}

\begin{grammarlens}[title={Grammar Lens}]
Keep the target beside its familiar subject and noun. The frame remains \textbf{\pa{ਮੈਂ} \pa{ਕੰਮ} \pa{ਕੀਤਾ}}, with \textbf{\pa{ਕੀਤਾ}} as the bounded completed form of familiar \textbf{\pa{ਕਰਨਾ}}.
\end{grammarlens}

\subsection*{Your turn}
\begin{itemize}
\raggedright
  \item \emph{Hear:} \emph{main kamm kiitaa}; YOU SAY: \textbf{\pa{ਮੈਂ} \pa{ਕੰਮ} \pa{ਕੀਤਾ}}
  \item \emph{Read it:} \textbf{\pa{ਮੈਂ} \pa{ਕੰਮ} \pa{ਕੀਤਾ}}
  \item \emph{Write it:} cover and write \textbf{\pa{ਮੈਂ} \pa{ਕੰਮ} \pa{ਕੀਤਾ}}
\end{itemize}

\begin{tcolorbox}[breakable,colback=teal!4,colframe=teal!35!black,title={Before you move on}]
Say the statement once more without looking.
\end{tcolorbox}
\section[\texorpdfstring{\pa{ਕਰਨਾ} · \pa{ਮੈਂ} \pa{ਕੰਮ} \pa{ਕੀਤਾ} (karnaa; main …)}{ਕਰਨਾ · ਮੈਂ ਕੰਮ ਕੀਤਾ (karnaa; main …)}]{Contrast action name and completed frame}

\label{lesson:PA-C190-infinitive-perfective-contrast}



\begin{quote}
Prepare the action name, then the completed work statement.
\end{quote}

\begin{grammarlens}[title={Grammar Lens: infinitive versus completed action}]
\begin{itemize}
\raggedright
  \item \textbf{\pa{ਕਰਨਾ}} names the action “to do.”
  \item \textbf{\pa{ਮੈਂ} \pa{ਕੰਮ} \pa{ਕੀਤਾ}} gives one completed masculine-singular frame.
\end{itemize}
\end{grammarlens}

\subsection*{Your turn}
\begin{itemize}
\raggedright
  \item \emph{Say it:} \textbf{\pa{ਕਰਨਾ}} for the action name
  \item \emph{Say it:} \textbf{\pa{ਮੈਂ} \pa{ਕੰਮ} \pa{ਕੀਤਾ}} for the completed action
  \item \emph{Write it:} \textbf{\pa{ਕਰਨਾ} · \pa{ਕੀਤਾ}}
\end{itemize}

\begin{tcolorbox}[breakable,colback=teal!4,colframe=teal!35!black,title={Before you move on}]
Cover the models and produce action name, then completed frame.
\end{tcolorbox}
\section[\texorpdfstring{\pa{ਕਰਨਾ} · \pa{ਮੈਂ} \pa{ਕੰਮ} \pa{ਕੀਤਾ} (karnaa; main …)}{ਕਰਨਾ · ਮੈਂ ਕੰਮ ਕੀਤਾ (karnaa; main …)}]{A bounded completed-action payoff}

\label{lesson:PA-C190-payoff}



\begin{quote}
Recall the infinitive and target, then prepare the complete frame.
\end{quote}

\subsection*{Your turn: four skills}
\begin{itemize}
\raggedright
  \item \emph{Hear:} \emph{karnaa; main kamm kiitaa}; YOU SAY: \textbf{\pa{ਕਰਨਾ}; \pa{ਮੈਂ} \pa{ਕੰਮ} \pa{ਕੀਤਾ}}
  \item \emph{Read it:} \textbf{\pa{ਕਰਨਾ} · \pa{ਮੈਂ} \pa{ਕੰਮ} \pa{ਕੀਤਾ}}
  \item \emph{Write it:} cover the line and write \textbf{\pa{ਕਰਨਾ} · \pa{ਕੀਤਾ}}
\end{itemize}

\begin{tcolorbox}[breakable,colback=teal!4,colframe=teal!35!black,title={Before you move on}]
Produce the action name and completed frame once more.

This is the thirteenth prerequisite tranche exposed by Punjabi A2 mock drafting. It does not unblock Punjabi A2 mock 1 by itself. It is not an A2 mock, not A2 readiness or attainment, and not external or human validation.
\end{tcolorbox}
